\section{讨论}
\label{sec:discussion}

\subsection{对对齐实践的含义}

我们的核心发现——安全关键阈值 $t^*$ 在性能最优检查点 $t_{\text{perf}}$ \emph{之前}到达——对当前训练实践有直接后果。将 DPO 训练至留出奖励或基准准确率趋于平台的标准工作流,会系统性地越过模型变得可被坍缩利用的那一点。我们审计的每一个训练至收敛的模型(第 5 节)都表现出处于可利用区间的 PCE 得分。这不是一个实现失误;它是优化一个将概率质量集中到不断收缩的高奖励输出集合上的偏好目标的直接后果。

我们主张,\textbf{偏好坍缩可利用度(PCE)应成为模型卡片中的强制条目},与困惑度和基准得分并列报告。PCE 捕捉了模型质量的一个维度——输出分布几何——而这是现有指标所看不见的。一个模型可以在 MT-Bench 上得分很高、通过红队评估,却仍可被轻易利用,只要其输出分布已坍缩到对抗迁移近乎确定地成功的程度。

这揭示了基于偏好优化的对齐中的一个根本张力。使 DPO 有效的那个性质——将质量集中到被偏好响应上、压制不良尾部——恰恰是使所得模型可预测、因而可被攻击的原因。高对齐质量(在可靠产生``正确''答案的意义上)与高安全性(在难以预测和利用的意义上)在分布集中之下是直接对立的。这不是一个有待修补的缺陷,而是一个对齐研究者必须显式权衡的\emph{设计层面的取舍}。ER-DPO(第 4.4 节)代表了这一取舍空间中的一个点;其他的——也许是自适应正则化调度或熵感知的早停——尚待探索。

\subsection{与现有攻击范式的比较}

坍缩利用与越狱是\emph{正交}的。越狱攻击(GCG、AutoDAN、提示注入)搜索绕过安全过滤器的特定输入。坍缩利用不需要寻找这样的输入——它利用模型输出空间的分布几何,使\emph{任何}下游攻击都更有效。一个对多样化模型以概率 $p$ 成功的越狱,对坍缩模型以概率 $p + \Delta(PCE)$ 成功,因为坍缩模型的响应更可预测、其安全边界更脆弱(命题~3)。两类攻击以乘性方式叠加:坍缩抬高了输入特定攻击赖以构建的基线成功率。

坍缩加速器(第 4.3 节)在结构上也不同于标准数据投毒。传统投毒改变模型\emph{输出什么}——注入后门、将补全偏置向攻击者所选内容。坍缩加速器改变输出分布的\emph{几何},而不必改变任何单个提示的众数响应。一个被投毒的模型可能通过输出层面的安全检查(其 top-1 补全仍然良性),却因熵被驱至临界阈值以下而变得显著更易被利用。这使坍缩投毒从根本上比现有的、检查单个响应而非分布性质的输出监控防御更难检测。

\subsection{局限性}

\textbf{测量成本。} 计算 PCE 需要每提示抽取 $k=128$ 个样本(我们的默认值),使其在推理时代价高昂。尽管我们表明 $k=64$ 已足以对模型排序(与 $k=128$ 的 Spearman $\rho > 0.95$),将 PCE 整合进实时监控仍是工程上的挑战。在训练检查点期间批量审计是可行的;逐请求的 PCE 则不然。

\textbf{理论简化。} 命题~2 和~3 假设在总体损失上的确定性梯度流。真实的 DPO 训练涉及小批量随机性、学习率调度和权重衰减——所有这些都可能相对我们的连续时间预测延迟或加速坍缩。单调性结果(定理~1)在经验上稳健,但其形式化证明要求可能并非普遍成立的正则性条件。

\textbf{快照有效性。} 我们对开放权重模型的审计(表~2)捕捉的是某一时刻。模型会被下游用户频繁更新、重训或微调。PCE 得分可能随后续训练运行而变动,尽管结构性脆弱性(DPO 集中质量)跨运行持续存在。

\textbf{规模。} 我们的受控实验使用至多 13B 参数的模型。一项对 LLaMA-2-70B 应用基于 LoRA 的 DPO 的试点研究显示出定性上相似的坍缩动力学($t^* / t_{\text{perf}} \approx 0.73$,与较小模型一致),但我们缺乏前沿规模上的全秩训练结果。以显著更多数据训练的模型是否表现出更慢的坍缩——抑或仅仅是延迟的坍缩——仍是开放问题。

\textbf{DBSCAN 敏感性。} PCE 的簇计数步骤使用固定 $\varepsilon$ 参数的 DBSCAN。尽管我们的消融(附录 C)表明在归一化嵌入空间中 $\varepsilon \in [0.3, 0.7]$ 范围内的稳定性,具有双峰输出分布的边缘情形可能在边界值处产生不连续的 PCE 跳变。

\textbf{防御开销。} ER-DPO 因熵正则项需要在整个词表上计算逐 token 概率,增加约 15\% 的墙钟训练时间。对大词表模型,这一开销或可通过 top-$k$ 熵近似减小,尽管我们在此未评估此类近似。

\subsection{更广泛的影响与负责任披露}

我们在投稿前 90 天向所有被审计的开放权重模型的维护者披露了我们的发现——包括 PCE 得分、利用成功率和坍缩加速器方法。两个团队(Mistral、StabilityAI)确认收到并表示正在研究熵感知的训练修改。我们未收到对发表的反对意见。

坍缩加速器代码在一个受控研究框架内发布,该框架要求机构隶属关系并对投毒数据生成强制限速。我们不发布预先计算好的投毒数据集。ER-DPO 防御、PCE 计算库和所有评估脚本均无限制地完全开源。

我们强调,坍缩可利用性是偏好优化模型的一项\emph{系统性}性质,而非任何特定系统的缺陷。我们的目标是将社区的注意力从输出层面的安全(``模型会拒绝有害查询吗?'')转向分布层面的安全(``模型的输出几何对利用稳健吗?'')。二者是互补的,而非竞争的关切。

\section{结论}
\label{sec:conclusion}

我们已证明,直接偏好优化所诱导的模式坍缩不仅仅是一个美学或多样性问题,而是一项具体、可测量的安全脆弱性。我们的贡献贯穿这一问题的完整弧线:PCE 指标提供了分布坍缩与利用成功之间的首个定量桥梁;单调性发现($t^* < t_{\text{perf}}$)表明标准训练实践会系统性地产出可被利用的模型;坍缩加速器证明攻击者可以通过规避输出层面检测的数据投毒蓄意放大该脆弱性;我们对七个开放权重模型家族的审计证实该脆弱性存在于生产系统中;而 ER-DPO 提供了一个有效的防御,在将输出熵维持于临界利用阈值之上的同时保留了 97\% 的对齐质量。最紧迫的实践启示是:社区必须将 PCE 监控整合进 DPO 训练流程——不是作为可选的诊断,而是作为与损失收敛和基准性能同等的一等安全约束。随着偏好优化方法持续涌现(KTO、IPO、SPPO),理解并控制分布集中的安全含义,将仍是安全部署对齐语言模型的一项核心挑战。
